\section{Split-base reflection, complete primary duality, and all mixed positive Euler forms}
\label{sec:split-base-mixed-duality}

This calculation starts with the original arithmetic multiplier
\[
 g(s)=2\xi(s)=s(s-1)\pi^{-s/2}\Gamma(s/2)\zeta(s)
\]
and a finite, nonempty symmetry-closed set $Z$ of its actual nontrivial
zeros, each with its actual full order $m_\lambda$. No auxiliary root
is admitted as an arithmetic zero. Retain
\[
 h(t)=\prod_{\lambda\in Z}(t-\lambda)^{m_\lambda},\quad
 d=\sum_\lambda m_\lambda,\quad E=\mathbb C[t]/(h),\quad
 A=M_t,\quad \varepsilon=j_h(h/g).
 \tag{SBD1}
\]
The source identities are CA.6--18, CA.26--31, AG37--38a,
DC.17--23, and TO.6--8. The source manifest pins their files.
The new calculation is the complete mixed-form classification SBD16--21
and its exact comparison with any original positive Gram in SBD22--24;
SBD2--15 construct and prove the particular support and duality maps
used in that classification. This result supplies no unproved purity
assertion and does not assert an actual offcritical zero.

\subsection{The absolute base and the reflected two-leg map}

For a complex vector space $X$ its split carrier is
\[
 G(X)=\{\tau\}\sqcup X^\bullet,\qquad e_X=0_X^\bullet\ne\tau.
 \tag{SBD2}
\]
Supported addition is vector addition, $\tau$ is the additive identity,
supported scalar multiplication is the original one, and $\tau$ is
absorbing for scalar multiplication. These structures lie over the
original pointed monoid blueprint
$\mathfrak b_\tau=\operatorname{Spec}\{\tau,1\}$ with
$\varnothing\equiv\tau$. For a linear map $L:X\to Y$ define
$G(L)(\tau)=\tau$ and $G(L)(x^\bullet)=(Lx)^\bullet$.
For an antilinear map the same formula is conjugate-semilinear.
Checking two supported arguments reduces every additive or scalar
identity to the corresponding identity for $L$; checking an argument
$\tau$ uses its defining rules. In particular $\ker L$ maps to $e_Y$,
whereas only external absence is mapped to external absence.

Let
\[
\begin{split}
 V&=\{\phi\in\mathcal S(\mathbb R):\phi(-x)=\phi(x),\
                  \phi(0)=0,\ \int_{\mathbb R}\phi=0\},\\
 \mathscr B&=\{F\in C^\infty(\mathbb R_{>0}):
        \sup_{x>0}x^b|D^kF(x)|<\infty
        \text{ for all }b\in\mathbb Z,k\geq0\},\quad D=-x\partial_x,\\
 \Theta\phi(x)&=2\sum_{n\geq1}\phi(nx),\qquad
 \widehat\phi(\nu)=\int_{\mathbb R}\phi(x)e^{-2\pi ix\nu}\,dx.
\end{split}
 \tag{SBD3}
\]
The original Fourier identities on $V$ are
$\widehat{\widehat\phi}=\phi$,
$\widehat{D\phi}=(1-D)\widehat\phi$ and
$x^{-1}\Theta\phi(1/x)=\Theta\widehat\phi(x)$.
For a retained $D$-invariant proper source $W\subseteq V$, put
\[
 W^\natural=\widehat{\overline W},\quad
 \mathfrak j_BF(x)=x^{-1}\overline{F(1/x)},\quad
 \mathfrak j_0(\phi,\psi)=(-\overline\psi,-\overline\phi).
 \tag{SBD4}
\]
The minus leg for $W$ is $\widehat W$, with action $1-D$;
the plus leg is $W$, with action $D$. Thus
$\mathfrak j_0$ maps the plus leg $W$ into the minus leg
$\widehat{W^\natural}=\overline W$, and the minus leg
$\widehat W$ into the plus leg $W^\natural$.
For a nonempty leg mask $S\subseteq\{+,-\}$ it has target mask
$S^{\mathrm{op}}$, obtained by exchanging the two leg names.
An absent source leg is omitted in this formula, not filled by a
supported zero. The target differential is the original differential
$\delta(\phi,\psi)=\Theta(\phi-\widehat\psi)$.
Direct substitution gives
\[
 \delta\mathfrak j_0=\mathfrak j_B\delta,\qquad
 \mathfrak j_0^2=1,\quad\mathfrak j_B^2=1,\qquad
 T\mathfrak j_0=\mathfrak j_0(1-T),\quad
 D\mathfrak j_B=\mathfrak j_B(1-D),
 \tag{SBD5}
\]
where $T=\operatorname{diag}(D,1-D)$ on the existing legs.
Indeed both sides of the first identity are
$\Theta(\widehat{\overline\phi}-\overline\psi)$, the
two squares cancel inversion or interchange, and the last two
identities follow by applying the displayed operators to a function.
Moreover $D W^\natural\subseteq W^\natural$, because
$D\widehat{\overline\phi}=\widehat{(1-D)\overline\phi}$.
Consequently this is a genuine map between the original properly
labelled cochain complexes. It is not an extra diagonal source on a
one-leg face. On a joint diagonal $(\phi,\widehat\phi)$ its
value is the joint diagonal
$(-\widehat{\overline\phi},-\overline\phi)$.

For $a>0$ retain $R_aF(x)=F(x/a)$ on the target and plus leg,
and $aR_{1/a}$ on the minus leg. This is a map from the
source label $W$ to the source label $R_aW$, since the
original Fourier substitution gives
$\widehat{R_a\phi}=aR_{1/a}\widehat\phi$.
In particular
$(R_aW)^\natural=aR_{1/a}W^\natural=R_{1/a}W^\natural$
as complex subspaces; the scalar $a$ is retained in the
actual maps below. No invariance of a proper $W$ under
$R_a$ is assumed. Substitution in SBD4 gives
\[
 \mathfrak j_B R_a=aR_{1/a}\mathfrak j_B,\qquad
 \mathfrak j_0\operatorname{diag}(R_a,aR_{1/a})
   =a\operatorname{diag}(R_{1/a},a^{-1}R_a)\mathfrak j_0.
 \tag{SBD6}
\]
These retain the original one-dimensional character $a$.
For an independent finite coefficient mask $I$, apply SBD4--6
coordinatewise on the face $X^I$. The face label remains $I$,
the leg label becomes $S^{\mathrm{op}}$, every killed vector
remains supported, and external $\tau$ remains external $\tau$.
Coordinatewise reflection commutes with every zero-insertion face
map $X^I\to X^J$, $I\subseteq J$, since it maps vector zero
to vector zero. Face insertion itself retains its original type;
in particular insertion from the empty face need not preserve
external absence.

\subsection{Exact arithmetic duality on every primary coordinate}

Write $\lambda^\dagger=1-\overline\lambda$ and
$f^\sharp(t)=\overline{f(1-\overline t)}$.
The original functional equations imply $g^\sharp=g$,
$h^\sharp=h$, and $\varepsilon^\sharp=\varepsilon$.
There are no real strip zeros in the original arithmetic input,
so $d$ is even. The full Chinese remainder coordinates are
\[
 E=\bigoplus_{\lambda\in Z}E_\lambda,\qquad
 E_\lambda=\mathbb C[x_\lambda]/(x_\lambda^{m_\lambda}),
 \quad A|_{E_\lambda}=\lambda+N_\lambda,\quad N_\lambda=M_{x_\lambda}.
 \tag{SBD7}
\]
Their idempotents are the original polynomials
$e_\lambda=j_h(h_\lambda r_\lambda)$, where
$h_\lambda=h/(t-\lambda)^{m_\lambda}$ and $r_\lambda$ is the
Taylor polynomial of $h_\lambda^{-1}$ at $\lambda$ through
degree $m_\lambda-1$. Taylor multiplication makes this
idempotent $1$ on its full primary factor and $0$ on every other
one. The two local actions of reflection are
\[
 e_\lambda^\sharp=e_{\lambda^\dagger},\qquad
 (x_\lambda^a e_\lambda)^\sharp
       =(-1)^a x_{\lambda^\dagger}^a e_{\lambda^\dagger}.
 \tag{SBD8}
\]
They follow by evaluating each primary jet; thus no nilpotent
coordinate has been replaced by its value.

Let $L_h$ be the top coefficient of the remainder modulo $h$.
For a remainder $P$, expansion of $P/h$ at infinity and its
partial fractions show that
$L_h(P)=\sum_\lambda\operatorname{Res}_\lambda(P/h)\,ds$.
Define
\[
 \mathcal R_g(f,u)=L_h(f^\sharp\varepsilon u)
   =\sum_\lambda\operatorname{Res}_{s=\lambda}
                         \frac{f^\sharp(s)u(s)}{g(s)}\,ds.
 \tag{SBD9}
\]
For completeness, $L_h(fu)$ is nondegenerate: if the remainder
$f$ has degree $r$ and nonzero leading coefficient $c$, then
$L_h(ft^{d-1-r})=c$. Both $\sharp$ and multiplication by the
unit $\varepsilon$ are invertible, so SBD9 is nondegenerate.
The top coefficient identity
$L_h(v^\sharp)=-\overline{L_h(v)}$ then proves
\[
 \mathcal R_g(f,u)=-\overline{\mathcal R_g(u,f)},\qquad
 \mathcal R_g(Af,u)+\mathcal R_g(f,Au)=\mathcal R_g(f,u).
 \tag{SBD10}
\]
The second equality follows directly from $(tf)^\sharp=(1-t)f^\sharp$.

Put $v_\mu(x)=g(\mu+x)/x^{m_\mu}$, retaining its entire
unit jet, and write $v_\mu(x)^{-1}=\sum_{r\geq0}b_{\mu,r}x^r$.
The complete primary matrix of SBD9 is
\[
 \mathcal R_g(e_\lambda x_\lambda^a,e_\mu x_\mu^b)
 =\begin{cases}
 (-1)^a b_{\mu,m_\mu-1-a-b},&\mu=\lambda^\dagger,
                                      \ a+b\leq m_\mu-1,\\
 0,&\text{otherwise}.
 \end{cases}
 \tag{SBD11}
\]
This follows by taking the coefficient of $x^{-1}$ in the
local fraction, and retains every coefficient of $g/h$ through
the required order. The anti-diagonal entries are
$(-1)^a/v_\mu(0)\ne0$. Thus duality exchanges exactly the
two reflected primary supports and pairs opposite levels of
their nilpotent filtrations.

More precisely, for $0\leq r_\lambda\leq m_\lambda$ put
$F_{\boldsymbol r}=\bigoplus_\lambda x_\lambda^{r_\lambda}E_\lambda$.
Its right orthogonal space for SBD9 is exactly
\[
 F_{\boldsymbol r}^{\perp_{\mathcal R_g}}
       =\bigoplus_{\mu\in Z}
                x_\mu^{m_\mu-r_{\mu^\dagger}}E_\mu.
 \tag{SBD12}
\]
In one paired block, multiplication by the unit $v_\mu^{-1}$
preserves every power ideal. Requiring its product with $u_\mu$
to have zero coefficients of degrees $0$ through
$m_\mu-1-r_{\mu^\dagger}$ is exactly divisibility by the
power displayed in SBD12. Those are all the coefficients tested
by pairing with $x^a$, $r_{\mu^\dagger}\leq a<m_\mu$.
This proves both inclusions, including the endpoint ideals.
On associated levels $a$ and $m_\mu-1-a$ the induced perfect
pairing has scalar $(-1)^a/v_\mu(0)$. SBD11 remains the full
pairing before passage to those levels. This is an exact
support-filtration correspondence, not an identification of a
supported-zero fibre with absence.

An offcritical orbit $\{\lambda,\lambda^\dagger\}$ has
equal full orders $m$. Each of its two primary summands is
totally isotropic for $\mathcal R_g$, and their cross matrix
$C$ in SBD11 is invertible. The Hermitian form $i\mathcal R_g$
on their sum has matrix
$\left(\begin{smallmatrix}0&iC\\-iC^*&0\end{smallmatrix}\right)$.
An invertible change of the second coordinate makes the cross
matrix the identity, and $(x,y)\mapsto(x+y,x-y)$ then exhibits
one positive and one negative copy of $\mathbb C^m$.
Consequently
\[
 \operatorname{sig}(i\mathcal R_g)|_{E_\lambda\oplus E_{\lambda^\dagger}}
       =(m,m,0).
 \tag{SBD13}
\]
This congruence proves inertia only; it does not replace the
original coordinates or unit in any subsequent formula.

The actual cohomological map is retained explicitly. With
$\phi_*=(4\pi^2x^4-6\pi x^2)e^{-\pi x^2}$, let
\[
 \begin{split}
 R_Z(u)&=\operatorname{rem}_h(\varepsilon u),\\
 R_{\mathrm{ref}}u&=S_h^{\mathscr B}\Theta(R_Z(u)(D)\phi_*),\qquad
 \sigma_Z(u)=[R_{\mathrm{ref}}u]\in Q[h(D)],\quad Q=\mathscr B/\Theta V.
 \end{split}
 \tag{SBD14}
\]
The original Euler inverse and jet theorem give
$\mathcal M R_{\mathrm{ref}}u=(g/h)R_Z(u)$ and the isomorphism
$\sigma_Z:E\to Q[h(D)]$. Because $g/h$ and $\varepsilon$
are sharp invariant and the sharp transform of a degree $<d$
polynomial still has degree $<d$, taking remainders commutes
with $\sharp$. Mellin transformation of $\mathfrak j_BF$
equals $(\mathcal MF)^\sharp$. Its injectivity on $\mathscr B$
(the original Mellin identification) therefore proves
\[
 \mathfrak j_B R_{\mathrm{ref}}=R_{\mathrm{ref}}\sharp,
 \qquad \mathfrak j_Q\sigma_Z=\sigma_Z\sharp,
 \qquad D\sigma_Z=\sigma_Z A.
 \tag{SBD15}
\]
Thus SBD9--13 and their coefficientwise versions have actual
isomorphisms to the arithmetic cohomology, with their complete
jets. The split lifts in SBD2 then put these on their original
support fibres over $\mathfrak b_\tau$. Nothing here asserts
that the finite section $R_{\mathrm{ref}}$ is $\mathbb C[t]$-linear:
its original nonzero boundary defect is retained in CA.31.

\subsection{Classification with arbitrary mixing on a coefficient face}

Fix a finite active coefficient mask $I$, $n=|I|>0$. On $E^I$
put $A_I=\operatorname{diag}_{i\in I}(A)$ and
\[
 \mathcal R_I(f,u)=\sum_{i\in I}\mathcal R_g(f_i,u_i).
 \tag{SBD16}
\]
Consider the vector space of all sesquilinear forms $B$ on
this fixed face that obey the exact complementary Euler identity
\[
 B(A_If,u)+B(f,A_Iu)=B(f,u).
 \tag{SBD17}
\]
We compute this entire vector space and its positive cone, allowing
all mixing between active coefficient copies. Nondegeneracy of
$\mathcal R_I$ gives a unique complex-linear operator $C$ with
$B(f,u)=\mathcal R_I(f,Cu)$. Substituting SBD10 in SBD17 gives
$\mathcal R_I(f,(CA_I-A_IC)u)=0$ for every $f,u$;
hence $CA_I=A_IC$. Write $C$ as an $I$-by-$I$ array of linear
operators on $E$. Each entry commutes with $A$. If $T$ is
such an entry, $T(P(A)1)=P(A)T(1)$, and every vector of $E$
is $P(A)1$. Thus $T=M_{T(1)}$. Conversely every multiplication
operator commutes with $A$. We have proved the exact bijection
\[
 \{B:\text{SBD17 holds}\}\ \longleftrightarrow\ M_I(E),
 \qquad B_C(f,u)=\sum_{i,j\in I}\mathcal R_g(f_i,c_{ij}u_j).
 \tag{SBD18}
\]
It retains the whole matrix of coefficient multipliers.
If $(C^\dagger)_{ij}=c_{ji}^\sharp$, moving multiplication
between the two arguments of SBD9 proves
$\overline{B_C(u,f)}=-B_{C^\dagger}(f,u)$. By uniqueness,
$B_C$ is Hermitian exactly when $C^\dagger=-C$.

The positive semidefinite cone in SBD18 is exactly
\[
 \boxed{\quad
 C=\sum_{\substack{\lambda\in Z\\\Re\lambda=1/2}}
           (e_\lambda j_hg')\,B_\lambda,
 \qquad B_\lambda\in M_I(\mathbb C),\quad B_\lambda=B_\lambda^*\succeq0.
 \quad}
 \tag{SBD19}
\]
Here multiplication by $e_\lambda j_hg'$ is applied to every
entry of the constant matrix $B_\lambda$. In particular all
offcritical primary blocks are zero, with no discarded
nilpotents in the assertion. We now prove both necessity and
sufficiency, including all degeneracies.

First a positive semidefinite Hermitian form satisfies
$|B(f,u)|^2\leq B(f,f)B(u,u)$. When $B(u,u)>0$, apply
nonnegativity to $f-z u$ and choose the complex scalar $z$
that minimizes the resulting quadratic. When $B(u,u)=0$,
nonnegativity of the same expression for all scalar sizes
forces $B(f,u)=0$. In particular a vector of zero quadratic
value is in the full radical.
For an offcritical $\lambda$, SBD11 and SBD18 show that
$B_C$ vanishes on $(E_\lambda)^I\times(E_\lambda)^I$.
Every vector in this space is therefore radical. Pairing with
$(E_{\lambda^\dagger})^I$ and using the nondegeneracy of
the cross matrix in SBD11 shows that every entry of $C$
on $E_\lambda$ vanishes. Repeating for its partner proves
the entire offcritical assertion, without assuming that $C$
is a scalar or invertible.

Now let $\lambda=1-\overline\lambda$. Distinct critical
primary summands are orthogonal by SBD11. Work on the single
summand $(E_\lambda)^I$ and put $N=A_I-\lambda I$.
Equation SBD17 becomes $B(Nf,u)+B(f,Nu)=0$.
Its radical is $N$-stable by this equality; on the quotient
the induced form is positive definite and the induced $N$
is nilpotent and skew-adjoint. It is zero: for real $r$,
differentiate the squared norm of the finite polynomial
$\exp(rN)v$ to obtain zero. If the largest $j$ with
$N^jv\ne0$ were positive, its squared norm would have
strictly positive leading coefficient
$\|N^jv\|^2/(j!)^2$ at degree $2j$, contradicting
constancy. This proves $N(E_\lambda)^I$ lies in the radical.
Consequently the form factors through the original constant
jet vector $f(\lambda)\in\mathbb C^I$ and has the unique form
$m_\lambda f(\lambda)^*B_\lambda u(\lambda)$ with
$B_\lambda$ positive semidefinite Hermitian; the original
positive integer $m_\lambda$ has not been removed.

Finally the elementary local identity
\[
 \frac{g'(\lambda+x)}{g(\lambda+x)}
        =\frac{m_\lambda}{x}+\frac{v_\lambda'(x)}{v_\lambda(x)}
 \tag{SBD20}
\]
gives
$\mathcal R_g(f,e_\lambda j_hg' u)
=m_\lambda\overline{f(\lambda)}u(\lambda)$ at a critical
centre. This proves that its unique multiplier is exactly
$e_\lambda j_hg'$, including the original unit. Uniqueness
in SBD18 now proves necessity of SBD19. Conversely SBD20
gives, for a multiplier in SBD19,
\[
 B_C(f,u)=\sum_{\Re\lambda=1/2}
          m_\lambda f(\lambda)^*B_\lambda u(\lambda),
 \tag{SBD21}
\]
which is Hermitian and positive semidefinite. This proves
sufficiency. Equivalently on the full local algebra the
multiplier is $m_\lambda v_\lambda(0)x_\lambda^{m_\lambda-1}
B_\lambda$; this is the actual jet of $g'$, not a unit
chosen to remove $v_\lambda(0)$. If $Z$ is wholly offcritical,
the cone consists of the zero form. A given $B_C$ is positive
definite on the full nonempty face precisely when every centre
is critical, every $m_\lambda=1$, and every $B_\lambda$ is
positive definite. Thus the cone contains a positive definite
form precisely when every centre is critical and every order
is one. These statements follow immediately from the fully
displayed radical in SBD21.

More explicitly, for the canonical reduction
$\pi_{\mathrm{red}}:E^I\to\bigoplus_{\lambda\in Z}\mathbb C^I$,
$f\mapsto(f(\lambda))_\lambda$, the radical is exactly
\[
 \operatorname{rad}B_C
 =\pi_{\mathrm{red}}^{-1}\left(
       \bigoplus_{\Re\lambda\ne1/2}\mathbb C^I
       \ \oplus\!
       \bigoplus_{\Re\lambda=1/2}\ker B_\lambda\right).
 \tag{SBD21a}
\]
Indeed the nonnegative sum SBD21 vanishes on a vector
exactly when each critical constant vector belongs to the
corresponding matrix kernel; the proved Cauchy--Schwarz
statement identifies this null space with the radical.
This proves that the form detects every nonzero vector in
the complete reduced arithmetic spectrum exactly when all
centres are critical and all $B_\lambda$ are positive
definite, regardless of the full orders $m_\lambda$.
The map $\pi_{\mathrm{red}}$ commutes with Euler action
through multiplication by the actual centres, and its split
lift sends the entire nilpotent ideal to the supported zero
of the reduced fibre, never to $\tau$.
Thus positivity that is faithful on the reduced spectrum
does not impose simplicity of zeros or semisimplicity on
the full primary algebra. The stronger positive-definite
full-algebra assertion in the preceding paragraph is a
different, explicitly computed condition. No assertion
about Deligne's weight theory or its semisimplicity follows
from conflating these two conditions.

All of this permits arbitrary mixing on the fixed active
face. Restricting a pairing along zero insertion from a
smaller coefficient face selects the corresponding principal
submatrices of each $B_\lambda$, which remain positive
semidefinite because their quadratic values are obtained by
inserting zero coordinates. This is a statement about
pullback of pairings. It does not falsely assert that an
arbitrary matrix mixing coefficients commutes, as an operator,
with every zero-insertion arrow. The original leg masks and
proper-source labels are independent indices retained by the
maps SBD4--6 and SBD14--15.

\subsection{The exact foundational obstruction measured by a source Gram}

Let $G\succ0$ be any one of the programme's original positive
source Gram matrices on the unchanged coordinates of $E^I$.
Do not alter its masses or its coordinate frame. Let $R_I$
be the matrix of SBD16 in these same coordinates. Define
\[
 C_G=R_I^{-1}G,\qquad
 \mathcal W_G=A_I^*G+GA_I-G.
 \tag{SBD22}
\]
The map $C_G$ is a completely determined complex-linear
map from the original metric realization to the original
residue-dual realization. It is invertible since both $R_I$
and $G$ are invertible. Inserting
$A_I^*R_I+R_IA_I=R_I$ and $G=R_IC_G$ gives exactly
\[
 \boxed{\quad \mathcal W_G=R_I(C_GA_I-A_IC_G).
       \quad}
 \tag{SBD23}
\]
This is an identity on the complete retained coefficient
space, not a bound. It identifies the additive metric
defect with the failure of this specific arithmetic
comparison map to be $\mathbb C[t]$-linear. Thus the
comparison is an exact module map if and only if the
original defect is zero. In that case SBD19 applies to
$G$, and its positive definiteness forces critical,
simple primary centres as just calculated. At an actual
offcritical packet the failure is already quantitatively
visible without discarding a single primary order: the
original vector $v=e_\lambda x_\lambda^{m_\lambda-1}$,
placed in any active coefficient, is a nonzero eigenvector,
so
\[
 v^*\mathcal W_Gv=(2\Re\lambda-1)v^*Gv,
 \qquad
 v^*R_I[C_G,A_I]v=(2\Re\lambda-1)v^*Gv.
 \tag{SBD24}
\]
This follows by substituting $A_Iv=\lambda v$ on both
terms, and $v^*Gv>0$ is its original norm.

The correspondence SBD23 gives the foundational meaning of
the norm defect already calculated in the programme: it is
the arithmetic-duality commutator of an explicitly determined
map. An invertible arithmetic unit cannot remove it and
leave a positive, exact complementary form on an offcritical
packet, by the complete cone computation SBD19. Allowing
mixed coefficient matrices does not alter that conclusion.
The construction has not shown this commutator vanishes,
nor that its tensor-growth contribution has the bound needed
by the programme. It has computed its exact domain,
codomain, map, allowed positive cone and full primary radical,
and connected them to the original theta cohomology and
the two split-support legs. These are the proved results
of this foundational calculation.
